\subsection{\texorpdfstring{Equicontinuous subsets of \(\mathcal{L}(\mathbf{E};\mathbf{F})\)}{Equicontinuous subsets of \textbackslash mathcal\{L\}(\textbackslash mathbf\{E\};\textbackslash mathbf\{F\})}}

Let \(\mathrm{E}\) and \(\mathrm{F}\) be two locally convex spaces. For a subset \(\mathrm{H}\) of \(\mathcal{L}(\mathrm{E};\mathrm{F})\) to be equicontinuous it is necessary and sufficient that it is equicontinuous at the point 0 in \(\mathrm{E}\) (I, p.~9, prop. 6); this implies that for every neighbourhood \(\mathrm{V}\) of 0 in \(\mathrm{F}\) , the set \(\bigcap_{u\in \mathrm{H}}u^{-1}(\mathrm{V})\) is a neighbourhood of 0 in \(\mathrm{E}\) ; or that for every continuous semi-norm

\(p\) on \(\mathrm{F}\) , the function \(\sup_{u\in \mathrm{H}}(p\circ u)\) is a continuous semi-norm on E. Moreover (I, p.~5),

  H is uniformly equicontinuous. We note that the convex balanced envelope of an equicontinuous subset is equicontinuous, since if p is a continuous semi-norm on F and \(\tilde{H}\) the convex balanced envelope of H, we have, for the \(u_{i}\) in H, the inequality \(p \circ (\sum_{i} \lambda_{i} u_{i}) \leqslant \sum_{i} |\lambda_{i}| \cdot (p \circ u_{i})\) , hence \(\sup_{u \in \tilde{H}} (p \circ u) = \sup_{u \in H} (p \circ u)\) .

Consequently, the family of equicontinuous subsets is a convex bornology on \(\mathcal{L}(\mathrm{E};\mathrm{F})\) (III, p.~2, def. 2).

PROPOSITION 4. --- Let E, F be two locally convex spaces, and F be Hausdorff. Let the space \(F^{E}\) of all mappings from E into F be assigned the topology of simple convergence. Then

\begin{enumerate}
\def\labelenumi{(\roman{enumi})}
\item
  The set of linear mappings from \(\mathbf{E}\) into \(\mathbf{F}\) is closed in \(\mathbf{F}^{\mathrm{E}}\) .
\item
  If \(\mathbf{H}\) is an equicontinuous subset of \(\mathcal{L}(\mathrm{E};\mathrm{F})\) , the closure \(\overline{\mathbf{H}}\) of \(\mathbf{H}\) in \(\mathrm{F}^{\mathrm{E}}\) is contained in \(\mathcal{L}(\mathrm{E};\mathrm{F})\) and is equicontinuous.
\end{enumerate}

We know that \(\overline{\mathbf{H}}\) is equicontinuous (GT, X, § 2, No.~3, prop. 6). It remains to prove the assertion (i). Let \(x, y\) be in E and \(\lambda, \mu\) in K, and let \(\mathbf{A}(x, y, \lambda, \mu)\) be the set of all \(u \in \mathbf{F}^{\mathrm{E}}\) such that

\[
u (\lambda x + \mu y) - \lambda u (x) - \mu u (y) = 0.
\]

This set is closed in \(F^{E}\) since the mapping \(u \mapsto u(x)\) from \(F^{E}\) into F is continuous for every \(x \in E\) and since \(F\) is Hausdorff. But the set of linear mappings from \(E\) into \(F\) is equal to

\[
\bigcap_ {x, y, \lambda , \mu} \mathrm{A} (x, y, \lambda , \mu).
\]

Thus this set is closed in \(F^{E}\) .

COROLLARY 1. --- For an equicontinuous subset H of \(\mathcal{L}(\mathrm{E};\mathrm{F})\) to be relatively compact in \(\mathcal{L}_s(\mathrm{E};\mathrm{F})\) , it is necessary and sufficient that for all \(x\in \mathrm{E}\) , the set \(\mathrm{H}(x)\) of all \(u(x)\) as \(u\) ranges over \(\mathrm{H}\) , is relatively compact in \(\mathrm{F}\) .

In fact, this condition is necessary and sufficient for \(\overline{\mathbf{H}}\) to be compact in \(\mathrm{F}^{\mathrm{E}}\) (GT, I, § 9, No.~5, cor.).

COROLLARY 2. --- Every equicontinuous subset of the dual E' of E is relatively compact for the weak topology \(\sigma(E', E)\) on E' (III, p.~14, Example 4).

For, if H is an equicontinuous subset of \(E'\) , \(\sup_{u\in H}|u|\) is a continuous semi-norm on E; in particular, for every \(x\in E\) , the set \(\mathrm{H}(x)\) is bounded, hence relatively compact in the field of scalars.

COROLLARY 3. --- In the strong dual \(E_{b}^{\prime}\) of a semi-normed space E, every closed ball is compact for the weak topology \(\sigma(\mathrm{E}^{\prime},\mathrm{E})\) .

This ball is also closed for \(\sigma(\mathrm{E}^{\prime},\mathrm{E})\) .

PROPOSITION 5. --- Let E and F be two locally convex spaces and let T be a total subset of E. The following uniform structures coincide on every equicontinuous subset H of \(\mathcal{L}(\mathrm{E};\mathrm{F})\) :

\begin{enumerate}
\def\labelenumi{\arabic{enumi})}
\item
  the uniform structure of simple convergence in T;
\item
  the uniform structure of simple convergence in E;
\item
  the uniform structure of convergence in the precompact subsets of E.
\end{enumerate}

We recall (III, p.~15, prop. 2) that the \(\mathfrak{S}\) -topology on \(\mathcal{L}(\mathrm{E};\mathrm{F})\) coincides with the \(\tilde{\mathfrak{S}}\) -topology, where \(\tilde{\mathfrak{S}}\) is the smallest bornology adapted to E and containing \(\mathfrak{S}\) .

In the statement of prop. 5, we can therefore replace the word « total » by « everywhere dense ». The proposition then follows from the general properties of equicontinuous sets (GT, X, § 2, No.~4, th. 1).

Examples. --- * 1) Let \(\mu\) be the Lebesgue measure on \(\mathbf{R}\) , and let E be the semi-normed space \(\mathcal{L}^p(\mu)\) ( \(1 \leqslant p < \infty\) ) (INT, IV). For every numerical function \(f\) and every real number \(h\) , let \(f_h\) be the function \(x \mapsto f(x - h)\) . Clearly the mapping \(f \mapsto f_h\) defines a linear isometry from E onto itself. If \(f\) is continuous and has compact support, then \(f_h\) converges to \(f\) uniformly, hence also in the mean of order \(p\) , as \(h\) tends to 0. Since the set \(\mathcal{K}(\mathbf{R})\) of all continuous functions with compact support is dense in E, and the set of linear isometries of E is equicontinuous, it follows from prop. 5 that for every \(f \in E\) , \(f_h\) converges in the mean of order \(p\) to \(f\) as \(h\) tends to 0.

For \(p = 1\) , consider the Fourier transform, which associates to each \(f \in \mathcal{L}^1(\mu)\) the function \(\hat{f}\) on \(\mathbf{R}\) defined by

\[
\hat {f} (y) = \int e ^ {- 2 i \pi x y} f (x) d \mu (x).
\]

The set of linear forms \(f \mapsto \hat{f}(y)\) is an equicontinuous subset of the dual of \(\mathcal{L}^1(\mu)\) . On the other hand, we know that the set \(T\) of all characteristic functions of closed bounded intervals is a total subset of \(\mathcal{L}^1(\mu)\) ; and we verify easily that for all \(f \in T\) , the Fourier transform \(\hat{f}\) is a continuous function tending to zero at infinity. We deduce that this is true for all \(f \in \mathcal{L}^1(\mu)\) (« Riemann-Lebesgue theorem »).

The relation \(\sup_{y\in \mathbb{R}}|\hat{f} (y)|\leqslant \| f\| _1\) shows that the map \(f\mapsto \hat{f}\) is a continuous map

from \(\mathcal{L}^1 (\mu)\) into the space \(\mathcal{B}(\mathbf{R})\) of all bounded functions on \(\mathbf{R}\) , with the structure of uniform convergence. Since \(\hat{f}\) is continuous for all \(f\in \mathrm{T}\) , it follows that \(\hat{f}\) is continuous for all \(f\in \mathcal{L}^1 (\mu)\) . The fact that \(\hat{f}\) tends to zero at infinity follows from the fact that the subspace \(\mathcal{C}_0(\mathbf{R})\) of all continuous functions tending to zero at infinity is closed in \(\mathcal{B}(\mathbf{R})\) .

\begin{enumerate}
\def\labelenumi{\arabic{enumi})}
\setcounter{enumi}{1}
\tightlist
\item
  Let E be the space of all continuous numerical functions on R endowed with the topology of compact convergence. Let K be a compact subset of R and let \((\mu_{n})\) be a sequence of measures on R with support in K. Suppose \(\|\mu_{n}\|\leqslant1\) for all n.~The set of the \(\mu_{n}\) is then an equicontinuous subset of \(E'\) . Therefore, if for every function \(f\in E\) , we have \(\lim_{n\to\infty}\mu_{n}(f)=0\) , the sequence of functions \(x\mapsto\int e^{itx}d\mu_{n}(t)\) converges to 0,
\end{enumerate}

uniformly on every compact subset of \(\mathbf{R}\) (since the set of functions \(t \mapsto e^{itx}\) , as \(x\) ranges over a compact subset of \(\mathbf{R}\) , is compact in E). *

COROLLARY. --- Suppose F is Hausdorff. Let H be an equicontinuous subset of \(\mathcal{L}(\mathrm{E};\mathrm{F})\) . If a filter \(\Phi\) on H converges simply to a mapping \(u_{0}\) from E into F, then \(u_{0}\) is a continuous linear mapping from E into F, and \(\Phi\) converges uniformly to \(u_{0}\) on every precompact subset of E.

The first assertion follows from prop. 4 (III, p.~16) and the second from prop. 5 (III, p.~17).

PROPOSITION 6. --- Let H be an equicontinuous subset of \(\mathcal{L}(\mathrm{E};\mathrm{F})\) . If F is metrizable and if there exists a countable total set in E, then the uniform structure on H of simple convergence in E is metrizable. If in addition, there exists a countable total set in F, then there exists a countable everywhere dense set in H (for the topology of uniform convergence on compact subsets of E).

Let \((a_{n})\) be a total sequence in E. Then the mapping \(u \mapsto (u(a_{n}))\) is an isomorphism from \(\mathcal{L}(\mathrm{E};\mathrm{F})\) , where \(\mathcal{L}(\mathrm{E};\mathrm{F})\) has the uniform structure of simple convergence on the set of the \(a_{n}\) , onto a uniform subspace of \(F^{N}\) . If F is metrizable (resp. metrizable and satisfies the first axiom of countability) then this is also true for \(F^{N}\) (GT, IX, § 2, No.~4, cor. 2 and § 2, No.~8, corollary), and the proposition follows from prop. 5 (III, p.~17).

COROLLARY 1. --- Let E be a locally convex metrizable space, and F a normed space. Suppose that E and F both satisfy the first axiom of countability. Then \(\mathcal{L}(\mathrm{E};\mathrm{F})\) is the union of a countable family of equicontinuous subsets and there exists a countable set in \(\mathcal{L}(\mathrm{E};\mathrm{F})\) which is dense for the topology of uniform convergence on precompact subsets of E.

Let B be the unit ball of F and \((\mathrm{V}_{n})\) a countable fundamental system of neighbourhoods of 0 in E. For every integer n, the set \(H_{n}\) of all \(u \in \mathcal{L}(\mathrm{E}; \mathrm{F})\) such that \(u(\mathrm{V}_{n}) \subset \mathrm{B}\) is equicontinuous and \(\mathcal{L}(\mathrm{E}; \mathrm{F})\) is the union of the \(H_{n}\) . The corollary then follows from prop. 6.

COROLLARY 2. --- Every closed ball in the dual E' of a normed space satisfying the first axiom of countability, is a compact metrizable space for the weak topology \(\sigma(\mathrm{E}',\mathrm{E})\) , and for this topology there exists a countable dense subset in E'.

This follows from prop. 6 and from III, p.~17, cor. 3.

